\documentclass[11pt, a4paper]{article}

% --- Packages ---
\usepackage[margin=2.5cm]{geometry}
\usepackage{booktabs}
\usepackage{array}
\usepackage{parskip}
\usepackage{hyperref}
\usepackage{amsmath}
\usepackage{enumitem}
\usepackage{microtype}

% --- Hyperref setup ---
\hypersetup{
    colorlinks=true,
    linkcolor=black,
    urlcolor=blue,
    citecolor=black
}

% --- Title ---
\title{\textbf{Quantitative Dynamic Macroeconomics}\\[0.4em]
\large Project Instructions}
\author{Konstantin Gantert}
\date{this version: April 2026}

% ---------------------------------------------------------------
\begin{document}
\maketitle
\tableofcontents
\bigskip

% ---------------------------------------------------------------
\section{Overview}

The final part of your coursework consists of a take-home project completed in
groups of up to three students. You will extend a baseline New Keynesian DSGE
model, introduce a new shock, and use the Dynare toolbox to simulate and analyze
monetary policy.

The project follows a deliberate two-stage logic. In the \textbf{first stage},
you work with the model as a quantitative object: you implement your extension
and shock in Dynare, understand the impulse responses, and analyze how different
configurations of a Taylor-type monetary policy rule affect the economy. In the
\textbf{second stage}, you connect your model and findings to a real-world
economic episode, using the model as a lens rather than a data-fitting exercise.

The final deliverable is a \textbf{policy report} of at most 4,000 words
(approximately 8--10 pages, excluding appendix and references), addressed to a
central bank audience with an economics background. The word budget should be
spent primarily on describing and interpreting your simulations and on building
economic intuition around your results. Model equations are welcome in the main
text where they directly support the argument, but the report is not a technical
derivation: detailed algebra belongs in the appendix.

\subsection*{Files Provided}

\begin{itemize}[noitemsep]
    \item Project Instructions (this document)
    \item Baseline Model Handout
    \item Baseline Model Code for Dynare
    \item Paper Assessment Form
\end{itemize}

% ---------------------------------------------------------------
\section{The Baseline Model}

Your starting point is a standard closed-economy New Keynesian DSGE model, fully
described in the Baseline Model Handout. The key features relevant for the
project are summarised here.

\textbf{Households} maximise lifetime utility over consumption and leisure
subject to a cash-in-advance (CIA) constraint, a budget constraint, and a
capital accumulation equation. Labor supply is flexible and wages are fully
flexible.

\textbf{Capital} is owned by households and accumulates subject to
\textbf{quadratic investment adjustment costs} (IAC):
\[
    K_{t+1} = (1-\delta)K_t
    + \left(1 - \frac{\varphi}{2}\left(\frac{X_t}{X_{t-1}}-1\right)^2\right)X_t.
\]
IAC introduces a wedge between the market value of installed capital and its
replacement cost. The shadow price of the capital law of motion, $\tilde{\Theta}_t$,
represents \textbf{Tobin's $Q$} --- the ratio of the market value of an
additional unit of installed capital to its replacement cost. At steady state
$\tilde{\Theta} = \tilde{\mu}$, where $\tilde{\mu}$ is the CIA multiplier, so
Tobin's $Q$ equals one and there is no incentive to invest beyond the
depreciation replacement level. Away from steady state, $\tilde{\Theta}_t >
\tilde{\mu}_t$ signals that installed capital is more valuable than new
investment goods and triggers additional investment. Lagged investment
$X_{t-1}$ is a predetermined state variable in the baseline.

\textbf{Firms} are split into a competitive final goods sector and a continuum
of monopolistically competitive intermediate goods firms producing with a
Cobb-Douglas technology. Nominal price rigidity follows the Rotemberg (1982)
quadratic price adjustment cost specification.

\textbf{Fiscal policy} is passive and the government budget is balanced period
by period:
\[
    \tau_t = \tau_C C_t + \tau_L w_t L_t + \tau_K r_t K_t \quad \forall\,t,
\]
where $\tau_C$, $\tau_L$, $\tau_K$ are fixed tax rates on consumption, labor
income, and capital income respectively. All tax revenues are redistributed 
to the households lump-sum. There are no government bonds and no
debt dynamics. Fiscal policy is not a policy instrument in this project.

\textbf{Monetary policy} follows a Taylor-type interest rate rule, discussed in
detail in Task~3. This is the exclusive policy instrument.

\textbf{Shocks.} The baseline contains a technology shock $Z_t$ and a monetary
policy shock $m_t$, both AR(1) processes. Each topic adds one further shock.

\textbf{Steady state.} All five topic extensions are designed so that the
steady state of the extended model is identical to that of the baseline. Only
the dynamics change. This simplifies implementation and allows you to verify
your extension by checking that Dynare's steady-state solver produces the same
values as the baseline.

% ---------------------------------------------------------------
\section{Topics}

Each group selects one of the five topics below. Each topic specifies (i) a
model extension, (ii) a new shock, and (iii) monetary policy questions to
investigate. Topics~1--4 are the primary topics. Topic~5 is a reserve topic
available if additional groups register.

The \textbf{real-world episode} is chosen by your group, not assigned. Each
topic provides a characterisation of the episode type that fits the model and
shock. Your group identifies a specific historical episode, motivates the
choice, and uses it as the real-world anchor for the policy analysis. The
instructor reviews the episode choice at the check-in milestone before you
proceed to writing the report.

Throughout all topics, \textbf{fiscal policy remains passive and the government
budget is balanced}. There are no monetary--fiscal interactions. The exclusive
policy instrument is the Taylor-type interest rate rule.

% ---------------------------------------------------------------
\subsection{Topic 1 --- Variable Capital Utilization and a Cost-Push Shock}

\textbf{Model extension.} Allow households to choose the utilization rate $u_t$
of the capital stock they own. The representative household owns the aggregate
capital stock $K_t$ and rents effective capital $u_t K_t$ to the market at a
rental rate $r_t$ per unit of effective capital. Each firm $i$ then demands its
share of effective capital services at the market rate $r_t$; from the firm's
perspective only the effective capital bundle matters, not its decomposition
into $u_t$ and $K_t(i)$ separately.

Higher utilization raises output but accelerates capital depreciation. The
depreciation function is:
\[
    \delta(u_t,\,u_{t-1})
    = \delta_0
    + \delta_1(u_t-1)
    + \frac{\delta_2}{2}(u_t-1)^2
    + \frac{\delta_3}{2}(u_t-u_{t-1})^2,
\]
where $\delta_0$ is the baseline depreciation rate, $\delta_1 > 0$ ensures a
positive marginal depreciation cost at steady state, $\delta_2 > 0$ governs the
curvature of the utilization cost, and $\delta_3 \geq 0$ introduces
\textbf{quadratic utilization adjustment costs} --- penalising rapid changes in
utilization. At $\delta_3 = 0$ the standard frictionless utilization model is
recovered; as $\delta_3 \to \infty$ utilization is frozen at its steady-state
value of one, recovering the baseline without utilization. Both are useful
limiting cases for verifying your implementation.

The capital law of motion becomes:
\[
    K_{t+1} = \bigl(1-\delta(u_t,u_{t-1})\bigr)K_t
    + \left(1-\frac{\varphi}{2}\left(\frac{X_t}{X_{t-1}}-1\right)^2\right)X_t.
\]
The effective capital input in the firm's production function is $u_t K_t(i)$:
\[
    Y_t(i) = \bigl(Z_t L_t(i)\bigr)^{1-\alpha}
    \bigl(u_t K_t(i)\bigr)^\alpha
    \left[1-\frac{\kappa}{2}\left(\frac{P_t(i)}{P_{t-1}(i)}-1\right)^{\!2}\right].
\]
The utilization adjustment cost makes $u_{t-1}$ a new predetermined state
variable. The household's FOC for $u_t$ is a genuine Euler equation linking
the marginal benefit of utilization today (higher rental income) to its
intertemporal cost (faster depreciation today and higher adjustment costs
tomorrow). The steady state is unaffected since $u = 1$ implies
$\delta(1,1) = \delta_0$ and all adjustment cost terms vanish.

\textbf{Shock.} A \textit{cost-push shock} --- an exogenous, persistent
increase in the price markup of intermediate goods firms, entering the Phillips
curve directly. Formally, the elasticity of substitution $\rho_t$ follows an
AR(1) process around its steady-state value $\bar{\rho}$. This shock raises
inflation without an accompanying output boom, breaking the divine coincidence
and creating a genuine trade-off between inflation and output stabilization.
The interaction with utilization adjustment costs is key: firms cannot
immediately offset falling output by ramping up capital utilization --- the
adjustment cost constraint smooths the utilization response, making inflationary
dynamics more persistent.

\textbf{Monetary policy focus.} How does utilization rigidity ($\delta_3$)
change the optimal monetary policy response to a cost-push shock? Does a more
aggressive inflation response ($\theta_\pi$) help when utilization adjusts
slowly, or does it cause excessive output volatility? How does interest rate
smoothing ($\theta_R$) interact with the hump-shaped utilization response?

\textbf{Episode type.} A recessionary episode in which inflationary pressure
from supply-side cost increases played a central role, capacity utilization was
a visible margin of adjustment, and the central bank faced a clear trade-off
between fighting inflation and supporting output. Episodes in which a persistent
cost shock displaced inflation above target for several years are particularly
suitable.

% ---------------------------------------------------------------
\subsection{Topic 2 --- Consumption Habit Formation and a Discount Rate Shock}

\textbf{Model extension.} Introduce external habit formation in household
consumption preferences. The household utility function becomes:
\[
    U_t = \frac{(C_t - \varphi_C C_{t-1})^{1-\sigma}-1}
               {(1-\sigma)(1-\varphi_C)^{1-\sigma}}
          - \frac{\phi}{1+\gamma}L_t^{1+\gamma},
\]
where $\varphi_C \in [0,1)$ governs the strength of habits and $C_{t-1}$ is
lagged aggregate consumption, taken as given by the individual household. The
utility function is scaled by $(1-\varphi_C)^{1-\sigma}$ in the denominator so
that the marginal utility of consumption at steady state equals $\bar{C}^{-\sigma}$
--- identical to the baseline. The steady state is therefore unchanged; only
the dynamics differ.

At $\varphi_C = 0$ the scaling factor equals one and the baseline utility
function is recovered exactly. Habit formation introduces inertia in aggregate
consumption: households adjust spending slowly because they are anchored to
past consumption levels, making aggregate demand more persistent following
shocks. Lagged aggregate consumption $C_{t-1}$ becomes a new predetermined
state variable.

\textbf{Shock.} A \textit{discount rate shock} --- an exogenous, persistent
increase in the household discount factor $\beta_t$, following an AR(1) process
around its steady-state value $\bar{\beta}$. A higher discount factor makes
households more patient: they save more, defer consumption, and reduce aggregate
demand. This generates a demand-driven recession with disinflationary pressure.
The interaction with habit formation makes the contraction more persistent, as
the falling habit stock continues to depress consumption even after the shock
itself dissipates.

\textbf{Monetary policy focus.} How should the Taylor rule respond to a
persistent demand-driven recession? Analyze the role of the output response
coefficient ($\theta_Y$). Compare a strictly inflation-targeting rule to one
that also responds to output. How does the inertia introduced by habit formation
affect the optimal speed and aggressiveness of monetary easing?

\textbf{Episode type.} A recessionary episode driven primarily by a persistent
contraction in domestic demand, where monetary policy faced the challenge of
stimulating an economy with deeply entrenched pessimism, precautionary saving,
or household deleveraging. Episodes with prolonged below-target inflation
alongside weak output are particularly suitable.

% ---------------------------------------------------------------
\subsection{Topic 3 --- Labor Adjustment Costs and an Oil Supply Shock}

\textbf{Model extension.} Introduce quadratic labor adjustment costs (LAC) at
the firm level. Rapidly changing the workforce is costly --- reflecting
recruitment, training, and severance costs. This is captured as a resource loss
in the production technology:
\[
    Y_t(i) = \bigl(Z_t L_t(i)\bigr)^{1-\alpha-\nu}
    K_t(i)^\alpha\, O_t^\nu
    \left[1 - \frac{\kappa}{2}
    \left(\frac{P_t(i)}{P_{t-1}(i)} - 1\right)^{\!2}
    -\frac{\varphi_L}{2}
    \left(\frac{L_t(i)}{L_{t-1}(i)}-1\right)^{\!2}\right],
\]
where $\varphi_L \geq 0$ governs the magnitude of adjustment costs. At
$\varphi_L = 0$ the frictionless baseline labor market is recovered. LAC
introduces inertia in firms' labor demand: firms adjust employment gradually
rather than immediately, making hours worked more persistent following shocks.
Lagged aggregate labor $L_{t-1}$ becomes a new predetermined state variable.
The steady state is unaffected since adjustment costs are zero whenever
$L_t = L_{t-1}$.

\textbf{Shock.} An \textit{oil supply shock} --- an exogenous, persistent
contraction in the supply of oil, which enters as a third factor of production:
\[
    Y_t(i) = \bigl(Z_t L_t(i)\bigr)^{1-\alpha-\nu}
    K_t(i)^\alpha\, O_t^\nu,
\]
where $O_t$ is exogenous oil supply following an AR(1) process around its
steady-state value $\bar{O}$, $\nu > 0$ is the oil share in production, and
constant returns to scale are maintained ($\alpha+\nu < 1$). A negative oil
shock raises marginal costs and reduces output, generating stagflationary
dynamics. The interaction with LAC is key: firms want to shed labor in response
to higher costs, but adjustment frictions prevent rapid downsizing, keeping
marginal costs elevated for longer and amplifying the persistence of
inflationary pressure.

\textbf{Monetary policy focus.} How does labor market inertia change the
optimal monetary policy response to a stagflationary supply shock? Should the
central bank respond aggressively to inflationary pressure, or does this deepen
the output contraction given slow labor market adjustment? How do the oil share
$\nu$ and the adjustment cost parameter $\varphi_L$ affect the trade-offs
facing the central bank?

\textbf{Episode type.} A recessionary episode driven by a persistent
contraction in energy or commodity supply, where labor market rigidities
amplified the persistence of the downturn and monetary policy faced a prolonged
stagflationary trade-off. Episodes in which a central bank had to choose
between fighting inflation and supporting employment in the aftermath of an
energy shock are particularly suitable.

% ---------------------------------------------------------------
\subsection{Topic 4 --- Durable Consumption Goods and a Labor Supply Shock}

\textbf{Model extension.} Split household consumption into non-durable
consumption $C_t$ and a stock of durable goods $D_t$. The household utility
function becomes:
\[
    U_t = \frac{(C_t^\psi D_t^{1-\psi})^{1-\sigma}-1}
               {(1-\sigma)\,\psi^{1-\sigma}\,\Omega^{(1-\psi)(1-\sigma)}}
          - \frac{\phi}{1+\gamma}L_t^{1+\gamma},
\]
where $\psi \in (0,1)$ governs the expenditure share on non-durables and
$\Omega$ is a scaling constant defined below. The durable stock evolves
according to:
\[
    D_{t+1} = (1-\delta_D)D_t + X_{D,t},
\]
where $X_{D,t}$ is expenditure on new durable goods and $\delta_D \in (0,1]$
is the durable depreciation rate. At $\delta_D = 1$ durables fully depreciate
within one period, recovering the standard non-durable model --- a useful
special case for verifying the implementation. The resource constraint expands
to:
\[
    C_t + X_{D,t} + X_t = Y_t,
\]
where $X_t$ remains capital investment. The durable stock $D_t$ becomes a new
predetermined state variable with a structure directly analogous to capital
accumulation.

\textit{Hint on the scaling factor.} Without the denominator scaling, the
marginal utility of non-durable consumption at steady state would depend on
the steady-state durable-to-consumption ratio $\bar{D}/\bar{C}$, shifting the
steady state away from the baseline. To preserve steady-state invariance,
define the scaling constant:
\[
    \Omega \equiv \frac{\bar{D}}{\bar{C}}
    = \frac{1-\psi}{\psi}
      \cdot\frac{1-\beta(1-\delta_D)}{\delta_D}.
\]
This is the steady-state ratio of durable stock to non-durable consumption
implied by the household's intratemporal optimality condition between $C_t$ and
$D_t$. With this scaling, the marginal utility of $C_t$ at steady state reduces
to $\bar{C}^{-\sigma}$, identical to the baseline, so the steady state is
unchanged. Compute $\Omega$ once from the calibrated parameter values and treat
it as a fixed constant when implementing the model in Dynare.

\textbf{Shock.} A \textit{labor supply shock} --- an exogenous, persistent
increase in the disutility of labor, modeled as a shock to $\phi_t$ in the
household utility function, following an AR(1) process around its steady-state
value $\bar{\phi}$. This shock reduces households' willingness to supply labor
at any given wage, pushing up nominal wages, raising firms' marginal costs, and
reducing output. It operates simultaneously as a supply shock --- generating
inflationary pressure through higher marginal costs --- and as an income shock
--- reducing household labor income. The interaction with durable goods is
central: when labor income falls persistently, households sharply curtail
durable purchases while attempting to smooth non-durable consumption, generating
an asymmetric sectoral response that is well-documented empirically.

\textbf{Monetary policy focus.} Durable goods demand is highly sensitive to
interest rates because durables are investment-like: households compare the user
cost of the durable stock to the prevailing interest rate. A more aggressive
inflation response ($\theta_\pi$) stabilizes inflation but causes excessive
volatility in durable expenditure, potentially deepening the recession in the
durable sector. How does the presence of durable goods change the optimal Taylor
rule relative to a model without them? Does responding to output ($\theta_Y >
0$) help stabilize both sectors simultaneously? How does the durable share
$\psi$ affect the transmission of monetary policy?

\textbf{Episode type.} A recessionary episode in which a persistent disruption
to labor supply or household labor income led to a sharp and prolonged
contraction in durable goods spending, and monetary policy faced the challenge
of stabilizing both inflation and a highly interest-rate-sensitive component of
aggregate demand. Episodes with a visible collapse in housing, automobile, or
consumer durable purchases alongside labor market deterioration are particularly
suitable.

% ---------------------------------------------------------------
\subsection{Topic 5 --- Price Indexation and an Investment-Specific Technology
Shock}

\textbf{Model extension.} Add partial price indexation to past inflation by
modifying the quadratic price adjustment cost of intermediate goods firms. The
standard Rotemberg cost penalises deviations of the gross inflation rate from
one. With indexation, the reference point shifts so that firms are instead
penalised for deviating from last period's inflation rate:
\[
    \frac{\kappa}{2}
    \left(\frac{P_t(i)}{P_{t-1}(i)(1+\pi_{t-1})^\chi}-1\right)^{\!2} Y_t,
\]
where $\chi \in [0,1]$ is the degree of indexation. At $\chi = 0$ the standard
Rotemberg cost and the baseline Phillips curve are recovered exactly. This
modification enters the firm's profit maximisation problem directly: re-deriving
the FOC for $P_t(i)$ with the modified cost yields a hybrid Phillips curve in
which lagged inflation $\pi_{t-1}$ appears as a structural term with a clear
microeconomic interpretation --- it governs how strongly firms anchor their
reference price to past inflation. Lagged inflation $\pi_{t-1}$ becomes a new
predetermined state variable. The steady state is unaffected since at $\pi_t =
\pi_{t-1} = \pi^*$ the indexation term vanishes.

Indexation introduces \textbf{inflation persistence}: once inflation moves away
from target it returns more slowly, and the output cost of disinflation --- the
sacrifice ratio --- rises with $\chi$.

\textbf{Shock.} An \textit{investment-specific technology (IST) shock} --- an
exogenous, persistent improvement in the efficiency with which final goods are
transformed into capital goods. The capital law of motion becomes:
\[
    K_{t+1} = (1-\delta)K_t + A_t
    \left(1-\frac{\varphi}{2}\left(\frac{X_t}{X_{t-1}}-1\right)^{\!2}\right)X_t,
\]
where $A_t$ follows an AR(1) process around its steady-state value of one. A
positive IST shock reduces the effective price of investment, boosting capital
accumulation and output while putting downward pressure on marginal costs and
inflation. With price indexation, the resulting disinflationary pressure becomes
persistent: firms partially index their reference price downward to past low
inflation, pulling current inflation further below target.

\textbf{Monetary policy focus.} How does the degree of price indexation $\chi$
affect the central bank's ability to stabilize inflation following a supply-side
displacement? Analyze how the sacrifice ratio changes with $\chi$. Should the
central bank respond more or less aggressively to inflation when indexation is
high? What does this imply for the optimal values of $\theta_\pi$ and $\theta_R$?

\textbf{Episode type.} A recessionary or recovery episode in which supply-side
dynamics --- particularly investment or productivity developments --- created
persistent disinflationary or deflationary pressure, and monetary policy faced
the challenge of preventing inflation from drifting too far below target.
Episodes in which a central bank underestimated inflation persistence or faced
a prolonged below-target inflation environment are particularly suitable.

% ---------------------------------------------------------------
\section{Tasks}

\subsection{Task 1 --- Study and Extend the Baseline Model}

Study the provided Dynare baseline code and the Baseline Model Handout
carefully. Identify the key equations, state variables, and parameters before
making any changes. Then implement the model extension specified for your topic.
You do not need to re-derive the entire model from scratch. Focus on:

\begin{itemize}[noitemsep]
    \item the additional first-order conditions introduced by your extension,
    \item any new state variables and the equations governing their evolution,
    \item a brief plain-language explanation of the economic mechanism your
          extension adds and why it matters for monetary policy transmission.
\end{itemize}

\textbf{Model summary table.} Include a table in the appendix of your report
listing all variables and all equations of your extended model. Completing this
table serves as a consistency check: a mismatch between the number of equations
and variables is the most common source of Dynare errors, and the table forces
you to verify that the model is internally complete before you simulate.

\textbf{Steady-state verification.} Before running stochastic simulations,
verify that your steady state is correctly implemented. Since all extensions are
designed to leave the steady state unchanged, Dynare's steady-state solver
should produce values identical to the baseline. Report the key steady-state
values in the calibration section of your report.

A working Dynare \texttt{.mod} file with the extension implemented, together
with a short explanatory note of approximately half a page, are required for the
instructor check-in (see Section~\ref{sec:milestones}).

\subsection{Task 2 --- Introduce the Shock and Analyze Baseline Dynamics}

Add the shock process specified for your topic to the Dynare model as an AR(1)
process:
\[
    \varepsilon_t = \rho_\varepsilon\,\varepsilon_{t-1} + \eta_t,
    \qquad \eta_t \sim \mathcal{N}(0,\,\sigma_\varepsilon^2).
\]
Calibrate the shock persistence $\rho_\varepsilon$ and standard deviation
$\sigma_\varepsilon$ to produce a recession of plausible size --- approximately
a 2--3\% peak decline in output. Analyze the impulse response functions (IRFs)
of the model and address the following:

\begin{itemize}[noitemsep]
    \item What are the key transmission channels from the shock to output and
          inflation?
    \item What happens to output, inflation, hours worked, investment, and the
          nominal interest rate?
    \item Do the IRFs tell a coherent economic story consistent with the
          mechanism your extension introduces? If not, revisit parameters.
\end{itemize}

\textbf{Model statistics.} Report a table of model moments: at minimum the
standard deviations of output, inflation, hours worked, investment, and the
nominal interest rate, as well as their contemporaneous correlations with
output. Dynare computes these automatically via the \texttt{stoch\_simul}
command.

For the moments table, activate \textbf{all three shocks simultaneously} ---
your topic shock, the baseline technology shock $Z_t$, and the baseline monetary
policy shock $m_t$. A single shock is typically insufficient to match the full
covariance structure of the data; combining shocks gives a richer and more
realistic picture of the model's business cycle properties.

As a benchmark, compare your model moments against the following stylized
business cycle facts for advanced economies, drawn from Stock, J.H.\ and
Watson, M.W.\ (1999), ``Business Cycle Fluctuations in U.S.\ Macroeconomic
Time Series'', in J.B.\ Taylor and M.\ Woodford (eds.), \textit{Handbook of
Macroeconomics}, Vol.~1A, Elsevier: output volatility around 1.5\%, consumption
less volatile than output, investment two to three times more volatile than
output, hours roughly as volatile as output, inflation less volatile than output,
and positive co-movement of consumption, investment, and hours with output.

Do not merely report the numbers --- discuss what the pattern of volatilities
and co-movements implies economically. Where does your model match the empirical
benchmark well, and where does it fall short? How does each shock contribute to
the overall moments?

\textit{Hint: A useful way to understand each shock's contribution is to
activate the shocks one at a time first, observe how the moments change, and
then combine them. This also helps you calibrate the relative magnitudes of the
shock standard deviations.}

\subsection{Task 3 --- Monetary Policy Analysis}

This is the analytical core of the project. The monetary policy instrument is a
Taylor-type interest rate rule of the form:
\[
    \frac{1+r_{B,t}}{1+\bar{r}_B}
    =
    \left(\frac{1+r_{B,t-1}}{1+\bar{r}_B}\right)^{\!\theta_R}
    \!\left[
        \left(\frac{1+\pi_t}{1+\pi^*}\right)^{\!\theta_\pi}
        \left(\frac{Y_t}{\bar{Y}}\right)^{\!\theta_Y}
    \right]^{1-\theta_R}
    e^{m_t},
\]
where $r_{B,t}$ is the net nominal interest rate, $\pi^*$ is the inflation
target, $\bar{r}_B$ and $\bar{Y}$ are steady-state values, and $m_t$ is a
monetary policy shock. The coefficient $\theta_R \in [0,1)$ governs interest
rate smoothing, $\theta_\pi > 1$ ensures the Taylor principle is satisfied, and
$\theta_Y \geq 0$ controls the response to output fluctuations.

Conduct systematic policy experiments. At minimum:

\begin{itemize}[noitemsep]
    \item Vary the inflation response coefficient $\theta_\pi$ across at least
          three values and compare IRFs.
    \item Vary the output response coefficient $\theta_Y$ across at least three
          values.
    \item Consider at least one alternative rule specification motivated by your
          topic's monetary policy focus --- for example, targeting output growth
          instead of the output level, or varying the smoothing coefficient
          $\theta_R$.
\end{itemize}

Use Dynare \texttt{for}-loops to generate overlaid IRF plots across policy
configurations. For each experiment, provide economic interpretation. What
trade-off does the central bank face? Which rule configuration performs best,
and by what criterion?

\subsection{Task 4 --- Connect to the Real-World Episode}

Having understood your model's mechanics, identify a real-world economic episode
that fits the characterisation given for your topic and connect the analysis to
it. Your episode choice must be approved at the instructor check-in before you
proceed to writing the report.

This task does not require formal econometric calibration to country data.
However, as optional enrichment, you are encouraged to use publicly available
data --- such as the FRED database of the St.\ Louis Fed
(\url{https://fred.stlouisfed.org}) or Eurostat
(\url{https://ec.europa.eu/eurostat}) --- to calibrate steady-state targets
such as the labor share, the investment-to-output ratio, or the average
inflation rate to the country in your episode. Document any such sources
clearly.

At minimum, address the following:

\begin{itemize}[noitemsep]
    \item Describe the episode in one to two paragraphs: what happened, and why
          does your extension and shock provide a useful lens?
    \item Identify the monetary policy stance actually taken and map it
          approximately onto a specific Taylor rule configuration in your model.
    \item Compare historical policy with the counterfactual experiments from
          Task~3. Does your model suggest the central bank could have done
          better? Be honest about the limitations of the mapping between model
          and episode.
\end{itemize}

\subsection{Task 5 --- Write the Policy Report}

Write the final report following the template in Section~\ref{sec:template}.
The report is addressed to a central bank economist with a macroeconomics
background. Lead with economic intuition; use model output --- IRFs, moments
tables, and steady-state values --- as supporting evidence, not as ends in
themselves. The report must include at minimum: one baseline IRF figure, one
policy experiment figure with overlaid IRFs across at least three rule
configurations, and one model moments table.

% ---------------------------------------------------------------
\section{Milestones and Deadlines}
\label{sec:milestones}

\begin{center}
\begin{tabular}{p{3.5cm} p{7.5cm} p{2.7cm}}
\toprule
\textbf{Milestone} & \textbf{Description} & \textbf{Deadline} \\
\midrule
Topic selection
    & Top two topic preferences and group composition
    & End of Week 8 \\[0.5em]
Instructor check-in
    & Working \texttt{.mod} file with extension, short explanatory
      note, and proposed real-world episode with brief motivation
    & End of Week 11 \\[0.5em]
Peer feedback session
    & Draft covering Tasks~1--3 for in-class peer review
    & Start of Week 13 \\[0.5em]
Final report
    & Complete policy report and Dynare code
    & End of Week 14 \\
\bottomrule
\end{tabular}
\end{center}

\bigskip
The instructor check-in serves two purposes: verifying that the Dynare
implementation is correct before work on the report begins, and approving the
episode choice before it shapes the policy narrative.

% ---------------------------------------------------------------
\section{Policy Report Template}
\label{sec:template}

The template below is a guideline, not a rigid prescription. Approximate lengths
are given for orientation based on a maximum of 4,000 words, excluding appendix
and references.

\begin{enumerate}

    \item \textbf{Executive Summary} (max half page)

    \item \textbf{Introduction} (approx.\ one page)
    \begin{enumerate}[label=(\alph*)]
        \item The economic episode: what happened and why it matters
        \item What the model captures and what question you are asking
        \item Brief preview of your main policy finding
    \end{enumerate}

    \item \textbf{Model Overview} (approx.\ 1.5 pages)
    \begin{enumerate}[label=(\alph*)]
        \item Structure of the model in words and/or a diagram
        \item The model extension: what mechanism it adds and why it matters for
              your episode
        \item Calibration: key parameter choices with brief justification (and
              sources)
    \end{enumerate}

    \item \textbf{Shock and Baseline Dynamics} (approx.\ 1.5 pages)
    \begin{enumerate}[label=(\alph*)]
        \item The shock: what it represents, how it is parameterized, and what
              the baseline IRFs look like
        \item Model statistics: a moments table with economic interpretation ---
              what do the volatilities and co-movements reveal about the model's
              dynamics given your shock?
        \item Connection to the real-world episode: do the baseline simulations
              capture the salient features of the historical recession?
    \end{enumerate}

    \item \textbf{Monetary Policy Experiments} (approx.\ 2 pages)
    \begin{enumerate}[label=(\alph*)]
        \item Systematic variation of the Taylor rule --- what you change and why
        \item Overlaid IRF comparison across policy configurations
        \item Economic interpretation of the trade-offs
    \end{enumerate}

    \item \textbf{Policy Recommendation} (approx.\ one page)
    \begin{enumerate}[label=(\alph*)]
        \item What does your analysis suggest the central bank should have done?
        \item Caveats: where does the model fall short as a guide to the episode?
        \item One key broader lesson about monetary policy design
    \end{enumerate}

    \item \textbf{References}

    \item \textbf{Appendix}\\
    Model variable and equation summary table, Dynare code (\texttt{.mod} file),
    additional IRF plots, and any derivations.

\end{enumerate}

% ---------------------------------------------------------------
\section{Rules and Academic Integrity}

\textbf{AI use} follows the course AI policy (Level~2). AI tools may be used
for grammar, spelling, and coherence checking of completed drafts only. AI may
not be used to generate content, code, tables, or graphs. All Dynare code must
be written by the group. Document AI tool use transparently in a technology
statement appended to the report.

\textbf{Dynare code} is delivered as a separate file and must run without errors.

All standard rules on plagiarism and academic integrity apply in full.

\end{document}
